\section{Related Work}

The AI CUP task uses the four-part formulation described in the PromiseEval
shared-task overview: identify an ESG promise, determine whether it has
actionable evidence, assess the clarity of the promise--evidence pair, and
assign a verification horizon
\cite{chen2025promiseeval}. ML-Promise is the associated
multilingual corpus resource; we use its English, French, Japanese, and Korean
releases for external replication \cite{seki2025mlpromise}. This joint
verification problem differs from
ClimateBERT-NetZero, which detects corporate and national net-zero or reduction
targets and supports analysis of their ambition, rather than jointly assessing
promise status, evidence, evidence quality, and timing
\cite{schimanski2023climatebert}.

\begingroup\emergencystretch=1em
These dependent outputs form a compact instance of hierarchical text
classification, in which labels occupy a supplied structure and coherent
predictions respect its ancestor relations \cite{plaud2024revisiting}. The AI
CUP hierarchy comes from the task's official field logic: we neither discover
the hierarchy nor propose a new hierarchical architecture. Instead, we compare
decision rules over the same field-wise probabilities while preserving or
ignoring that given structure.
\par\endgroup

Yu et al. introduced path-constrained MicroF1 and MacroF1, under which a
correct node receives credit only when all its ancestors are also correct
\cite{yu2022constrained}. Ji et al. later named this family \emph{C-metrics}
\cite{ji2023hierarchical}. Our path-constrained weighted macro-F1 (C-wF1)
adapts Yu et al.'s principle to this task's fields and its official
weighted macro-F1; it is therefore not the original C-MacroF1.
Section~\ref{sec:metrics} defines both it and the set-overlap hierarchical
F1 (hF). Both were adopted post hoc and are reported as exploratory.

Plaud et al. observed from point estimates that constrained F1 variants did not
globally change method rankings in their experiments
\cite{plaud2024revisiting}. That observation and our post-hoc paired contrast
inference answer different questions: ranking a collection of models is not the
same estimand as a paired difference between decision rules on shared rows and
probabilities. Moreover, our hF changes both consistency treatment and
macro-versus-micro aggregation. Of the two post-hoc metrics, only C-wF1
isolates the treatment of children whose predicted ancestors do not
support them, while retaining
the official field aggregation.
